\chapter{التوزيع الثنائي}\label{ch:g11:binom}

كرّر التجربة ذات الجوابين نعم/لا عدة مرات، على نحو مستقل، ثم عُدّ
النجاحات: \omterm{def:g11:prob:rv}{فالتوزيع} الناتج --- \emph{\omterm{def:g11:prob:rv}{التوزيع} الثنائي} --- هو
أهم التوزيعات المتقطعة على الإطلاق. ويبني هذا الفصل هذا \omterm{def:g11:prob:rv}{التوزيع}
بالأشجار وعدّ المسارات؛ أما الصيغة المغلقة لأعداد المسارات (بالعوامل)
فتأتي مع أدوات العدّ في \cref{ch:g12:comb}، ويُعاد النظر في
\omterm{def:g11:prob:rv}{التوزيع} في \cref{ch:g12:randvar}.

\section{تجارب برنولي}

\begin{definition}[تجربة برنولي]\label{def:g11:binom:bernoulli}
\emph{تجربة برنولي}\index{تجربة برنولي} هي تجربة لها
إمكانيتان بالضبط: \emph{النجاح}، \omterm{def:g10:proba:distribution}{باحتمال} $p$، و
\emph{الفشل}، \omterm{def:g10:proba:distribution}{باحتمال} $1 - p$. ويقال عن \omterm{def:g11:prob:rv}{المتغير العشوائي} $X$ الذي يساوي
$1$ عند النجاح و $0$ عند الفشل إنه يتبع
\emph{\omterm{def:g11:prob:rv}{توزيع} برنولي} $\mathcal B(p)$؛ وعندئذٍ
\[
\E(X) = p,
\qquad
\V(X) = p(1 - p).
\]
\end{definition}

\begin{proof}[برهان الصيغتين]
$\E(X) = p \times 1 + (1-p) \times 0 = p$؛ وبما أن $X^2 = X$ (فكل من $0$
و $1$ مربع نفسه)، يكون $\E(X^2) = p$، إذن حسب
\cref{prop:g11:prob:konig}، $\V(X) = p - p^2 = p(1-p)$.
\end{proof}

\begin{definition}[تجارب مستقلة مكرَّرة]\label{def:g11:binom:repeated}
تكرار \omterm{def:g11:binom:bernoulli}{تجربة برنولي} $n$ مرة \emph{على نحو مستقل} يعني: أن
نتيجة كل تجربة لا تؤثر في الأخريات، وأن \omterm{def:g10:proba:distribution}{احتمال}
أي \omterm{def:g11:seq:sequence}{متتالية} كاملة من النتائج هو \emph{جداء}
الاحتمالات على طول المسار الموافق في \omterm{met:g10:proba:tree}{الشجرة} --- $p$ عند كل
نجاح، و $1 - p$ عند كل فشل.
\end{definition}

\begin{example}\label{ex:g11:binom:threetrials}
ثلاث تجارب مستقلة \omterm{def:g10:proba:distribution}{احتمال} النجاح فيها $p$. \omterm{def:g11:seq:sequence}{فالمتتالية} SFS
(نجاح، فشل، نجاح) احتمالها $p(1-p)p = p^2(1-p)$ --- وكذلك
\emph{كل} \omterm{def:g11:seq:sequence}{متتالية} فيها نجاحان بالضبط، مهما كانت
مواقعهما: إذ لا يهم إلا \emph{عدد} حروف S وحروف F.
\end{example}

\section{أعداد المسارات والمعاملات الثنائية}

\begin{definition}[المعامل الثنائي]\label{def:g11:binom:coefficient}
في \omterm{met:g10:proba:tree}{شجرة} $n$ تجربة مستقلة، \emph{المعامل
الثنائي}\index{معامل ثنائي} $\binom{n}{k}$ (ويُقرأ ``$n$
اختر $k$'') هو عدد المسارات التي تحوي $k$ نجاحًا بالضبط.
\end{definition}

\begin{example}
$\binom{3}{2} = 3$: وهي المسارات SSF و SFS و FSS. وبالمثل $\binom{3}{0} = 1$
(المسار FFF)، و $\binom{3}{1} = 3$ و $\binom{3}{3} = 1$. واصطلاحًا وبحكم \omterm{met:g10:proba:tree}{الشجرة}، يكون $\binom n0 = \binom nn = 1$ من أجل كل $n$.
\end{example}

\begin{omfigure}
\begin{tikzpicture}[font=\small]
  \coordinate (root) at (0,0);
  \node (S) at (1.5,1.5) {S};
  \node (F) at (1.5,-1.5) {F};
  \node (SS) at (3.0,2.25) {S};
  \node (SF) at (3.0,0.75) {F};
  \node (FS) at (3.0,-0.75) {S};
  \node (FF) at (3.0,-2.25) {F};
  \node (SSS) at (4.5,2.6) {S};
  \node (SSF) at (4.5,1.85) {F};
  \node (SFS) at (4.5,1.1) {S};
  \node (SFF) at (4.5,0.35) {F};
  \node (FSS) at (4.5,-0.35) {S};
  \node (FSF) at (4.5,-1.1) {F};
  \node (FFS) at (4.5,-1.85) {S};
  \node (FFF) at (4.5,-2.6) {F};
  \draw[omThm, thick] (root) -- (S);
  \draw[omThm, thick] (S) -- (SS);   \draw[omThm, thick] (SS) -- (SSF);
  \draw[omThm, thick] (S) -- (SF);   \draw[omThm, thick] (SF) -- (SFS);
  \draw[omThm, thick] (root) -- (F);
  \draw[omThm, thick] (F) -- (FS);   \draw[omThm, thick] (FS) -- (FSS);
  \draw[gray] (SS) -- (SSS);
  \draw[gray] (SF) -- (SFF);
  \draw[gray] (FS) -- (FSF);
  \draw[gray] (F) -- (FF);
  \draw[gray] (FF) -- (FFS);
  \draw[gray] (FF) -- (FFF);
  \node[gray, anchor=west] at (5.0,2.6) {$3$ نجاحات};
  \node[omThm, anchor=west] at (5.0,1.85) {$2$ نجاحان};
  \node[omThm, anchor=west] at (5.0,1.1) {$2$ نجاحان};
  \node[gray, anchor=west] at (5.0,0.35) {$1$ نجاح};
  \node[omThm, anchor=west] at (5.0,-0.35) {$2$ نجاحان};
  \node[gray, anchor=west] at (5.0,-1.1) {$1$ نجاح};
  \node[gray, anchor=west] at (5.0,-1.85) {$1$ نجاح};
  \node[gray, anchor=west] at (5.0,-2.6) {$0$ نجاح};
\end{tikzpicture}

{\small \omterm{met:g10:proba:tree}{شجرة} $n = 3$ تجارب: $\binom{3}{2} = 3$ مسارًا (بالأحمر) تحمل
نجاحين بالضبط، كل منها \omterm{def:g10:proba:distribution}{باحتمال} $p^2(1-p)$.}
\end{omfigure}

\begin{proposition}[قاعدة باسكال]\label{prop:g11:binom:pascal}
\index{مثلث باسكال}
من أجل $1 \leq k \leq n - 1$:
\[
\binom{n}{k} = \binom{n-1}{k-1} + \binom{n-1}{k}.
\]
\end{proposition}

\begin{proof}
رتّب مسارات \omterm{met:g10:proba:tree}{شجرة} $n$ تجربة ذات $k$ نجاحًا حسب
تجربتها \emph{الأخيرة}. فالمسارات المنتهية بنجاح تُنال من
مسار في التجارب $n-1$ الأولى فيه $k - 1$ نجاحًا: وعددها
$\binom{n-1}{k-1}$. والمسارات المنتهية بفشل تمدّد مسارًا فيه
$k$ نجاحًا بين التجارب $n-1$ الأولى: وعددها $\binom{n-1}{k}$.
وكل مسار من أحد النوعين بالضبط.
\end{proof}

وتولّد قاعدة باسكال المعاملات سطرًا بعد سطر --- فكل خانة هي
مجموع الخانتين اللتين فوقها:
\[
\begin{array}{ccccccccccc}
&&&&&1&&&&&\\
&&&&1&&1&&&&\\
&&&1&&2&&1&&&\\
&&1&&3&&3&&1&&\\
&1&&4&&6&&4&&1&\\
1&&5&&10&&10&&5&&1
\end{array}
\]

\begin{remark}
تُقام صيغة مغلقة، $\binom nk = \frac{n!}{k!(n-k)!}$، مع
نظرية منهجية في العدّ، في \cref{ch:g12:comb}. أما في
هذا المستوى، \omterm{prop:g11:binom:pascal}{فمثلث باسكال} يحسب كل معامل نحتاج إليه.
\end{remark}

\section{التوزيع الثنائي}

\begin{theorem}[التوزيع الثنائي]\label{thm:g11:binom:binomial}
ليعدّ $X$ النجاحات في $n$ \omterm{def:g11:binom:bernoulli}{تجربة برنولي} مستقلة ذات
\omterm{def:g11:stat:median}{الوسيط} $p$. عندئذٍ يتبع $X$ \emph{التوزيع
الثنائي}\index{توزيع ثنائي} $\mathcal B(n, p)$:
\[
\P(X = k) = \binom{n}{k}\, p^k (1-p)^{n-k},
\qquad k = 0, 1, \dots, n .
\]
\end{theorem}

\begin{proof}
\omterm{def:g11:prob:model}{الحادثة} $X = k$ هي مجموعة كل المسارات التي فيها $k$
نجاحًا بالضبط. ولكل مسار منها \omterm{def:g10:proba:distribution}{الاحتمال} $p^k(1-p)^{n-k}$: فالجداء
على طول المسار يحوي $k$ عاملًا مقداره $p$ و $n - k$ عاملًا مقداره $1-p$، \omterm{def:g10:coordgeom:system}{بترتيب}
ما (\cref{def:g11:binom:repeated}). وعدد هذه المسارات $\binom nk$
(\cref{def:g11:binom:coefficient})، واحتمالاتها تُجمع.
\end{proof}

\begin{example}\label{ex:g11:binom:quiz}
اختبار قصير فيه $5$ أسئلة مستقلة، لكل منها $4$ خيارات؛ ويجيب تلميذ
عشوائيًا، فيكون كل سؤال نجاحًا \omterm{def:g10:proba:distribution}{باحتمال} $p = \frac14$. وعدد
الأجوبة الصحيحة $X$ يتبع $\mathcal B\left(5, \frac14\right)$،
وباستعمال السطر $5$ من \omterm{prop:g11:binom:pascal}{مثلث باسكال}:
\[
\P(X = 2) = \binom52 \left(\frac14\right)^2\left(\frac34\right)^3
= 10 \times \frac{1}{16} \times \frac{27}{64}
= \frac{270}{1024} \approx 0.26 .
\]
ويستعمل \omterm{def:g10:proba:distribution}{احتمال} جواب صحيح واحد على الأقل \omterm{def:g10:proba:operations}{المتممة}:
$\P(X \geq 1) = 1 - \P(X = 0) = 1 - \left(\frac34\right)^5 \approx
0.76$.
\end{example}

\begin{proposition}[الأمل الرياضي والتباين]\label{prop:g11:binom:expectation}
إذا كان $X \sim \mathcal B(n, p)$:
\[
\E(X) = np,
\qquad
\V(X) = np(1-p).
\]
\end{proposition}

\begin{proof}[تبرير]
اكتب $X = X_1 + X_2 + \dots + X_n$، حيث يساوي $X_i$ العدد $1$ إذا نجحت
التجربة رقم $i$: فكل $X_i$ متغير برنولي أمله الرياضي
$p$ (\cref{def:g11:binom:bernoulli}). والمتوسطات تُجمع --- فجمع الإسهامات $n$
يعطي $\E(X) = np$. أما أن \emph{التباينات} تُجمع أيضًا من أجل
المتغيرات المستقلة فصحيح لكنه أدق: فصيغة \omterm{def:g11:prob:variance}{التباين}
\emph{مقبولة في هذا المستوى} ومبرهَن عليها في \cref{ch:g12:sums}.
\end{proof}

\begin{omfigure}
\begin{tikzpicture}
\begin{axis}[
  width=10cm, height=5cm,
  ybar, bar width=9pt,
  xlabel={$k$}, ylabel={$\P(X = k)$},
  xmin=-0.8, xmax=10.8, ymin=0, ymax=0.3,
  xtick={0,1,2,3,4,5,6,7,8,9,10},
  ytick={0.1, 0.2},
  tick label style={font=\small},
  label style={font=\small},
  title={$\mathcal B(10,\ 0.5)$}]
  \addplot[fill=omDef!40, draw=omDef] coordinates {
    (0,0.00098) (1,0.00977) (2,0.04395) (3,0.11719) (4,0.20508)
    (5,0.24609) (6,0.20508) (7,0.11719) (8,0.04395) (9,0.00977)
    (10,0.00098)};
\end{axis}
\end{tikzpicture}

{\small \omterm{def:g11:prob:rv}{التوزيع} $\mathcal B(10, 0.5)$ (عشر رميات لقطعة نقد غير مغشوشة):
مركزه $\E(X) = np = 5$، وهو متناظر، وكل \omterm{def:g10:proba:distribution}{الاحتمال} \omterm{def:g10:numbers:approx}{تقريبًا}
بين $2$ و $8$.}
\end{omfigure}

\begin{method}[التعرف على وضعية ثنائية]\label{met:g11:binom:recognize}
قبل أن تكتب $X \sim \mathcal B(n, p)$، تحقق من ثلاثة مكونات:
\emph{عدد ثابت} $n$ من التجارب، مقرَّر مسبقًا؛ ولكل تجربة
\emph{إمكانيتان} \omterm{def:g10:proba:distribution}{باحتمال} النجاح \emph{نفسه} $p$؛ والتجارب
\emph{مستقلة} (بالإرجاع، أو بأجهزة منفصلة).
أما السحب بدون إرجاع من مجتمع صغير \emph{فليس}
ثنائيًا --- إذ يتغير \omterm{def:g10:proba:distribution}{الاحتمال} عند كل سحبة
(\cref{exo:g11:prob:6}).
\end{method}

\section{المعاينة: هل الملاحظة مفاجئة؟}

يجيب \omterm{def:g11:prob:rv}{التوزيع} الثنائي عن سؤال عملي جدًا: \emph{إذا كان \omterm{def:g10:proba:distribution}{احتمال}
النجاح $p$ فعلًا، فما أعداد النجاحات
المعقولة؟}

\begin{example}\label{ex:g11:binom:decision}
يُفترض بآلة أن تنتج $10\%$ من القطع معيبة على الأكثر. وفي
دفعة من $10$ قطع، وُجدت $4$ معيبة. أهو سوء حظ أم عطب في الآلة؟ فإذا
كانت الآلة سليمة، تبع عدد المعيبات
$\mathcal B(10, 0.1)$، ويكون
\[
\P(X \geq 4) = 1 - \P(X \leq 3) \approx 1 - 0.987 = 0.013 :
\]
أي نحو فرصة واحدة من $80$. وملاحظة \omterm{def:g11:prob:model}{حادثة} بهذا القدر من عدم \omterm{def:g10:proba:distribution}{الاحتمال}
إشارة قوية --- \emph{فتُرفض} الفرضية القائلة إن الآلة ما زالت تعمل
عند $10\%$، مع تذكّر أن القرار قد يكون خاطئًا \omterm{def:g10:proba:distribution}{باحتمال}
نحو $0.013$.
\end{example}

\begin{method}[قاعدة قرار انطلاقًا من نموذج ثنائي]\label{met:g11:binom:decision}
للحكم على عدد ملاحَظ $k$ من النجاحات في مواجهة فرضية
$X \sim \mathcal B(n, p)$: احسب \omterm{def:g10:proba:distribution}{احتمال} الحصول، في ظل الفرضية،
على نتيجة \emph{لا تقل تطرفًا} عن $k$. فإذا كان ذلك \omterm{def:g10:proba:distribution}{الاحتمال}
صغيرًا جدًا (والاصطلاح الشائع: دون $5\%$)، فارفض الفرضية؛
وإلا فالملاحظة متوافقة معها. والعتبة اختيار،
لا مبرهنة --- فالإحصاء يكمّم الخطر، والمستعمل
يقبله.
\end{method}

\section{تمارين}

\begin{exercise}[$\star$]\label{exo:g11:binom:1}
مدّد \omterm{prop:g11:binom:pascal}{مثلث باسكال} إلى السطر $7$، وأعطِ قيم
$\binom62$ و $\binom{7}{3}$ و $\binom74$.
\end{exercise}

\begin{exercise}[$\star$]\label{exo:g11:binom:2}
يُرمى حجر نرد غير مغشوش $4$ مرات؛ ويعدّ $X$ عدد الأوجه الستة. برّر أن
$X \sim \mathcal B\left(4, \frac16\right)$ ثم احسب $\P(X = 0)$
و $\P(X = 1)$ و $\P(X \geq 2)$.
\end{exercise}

\begin{exercise}[$\star$]\label{exo:g11:binom:3}
أي مما يلي ثنائي؟ برّر.
\begin{enumerate}
  \item عدد الصور في $20$ رمية لقطعة نقد غير مغشوشة.
  \item عدد الآصات في $5$ أوراق تُوزَّع من رزمة واحدة.
  \item عدد الأيام الممطرة في الأسبوع المقبل، إذا كان كل يوم ممطرًا
        \omterm{def:g10:proba:distribution}{باحتمال} $0.3$ على نحو مستقل.
\end{enumerate}
\end{exercise}

\begin{exercise}[$\star$]\label{exo:g11:binom:4}
$X \sim \mathcal B(50, 0.2)$. أعطِ $\E(X)$ و $\V(X)$ و $\sigma(X)$.
\end{exercise}

\begin{exercise}[$\star\star$]\label{exo:g11:binom:5}
رامٍ يصيب الهدف \omterm{def:g10:proba:distribution}{باحتمال} $0.7$ في كل رمية،
على نحو مستقل. في $6$ رميات، احسب \omterm{def:g10:proba:distribution}{احتمال} $4$ إصابات بالضبط،
\omterm{def:g10:proba:distribution}{واحتمال} $5$ إصابات على الأقل.
\end{exercise}

\begin{exercise}[$\star\star$]\label{exo:g11:binom:6}
اختبار صواب/خطأ فيه $8$ أسئلة؛ ويخمّن تلميذ كل جواب.
فما \omterm{def:g10:proba:distribution}{احتمال} النجاح ($6$ أجوبة صحيحة على الأقل)؟
\end{exercise}

\begin{exercise}[$\star\star$]\label{exo:g11:binom:7}
كل علبة حبوب تُشترى تحوي التمثال A أو التمثال B، \omterm{def:g10:proba:distribution}{باحتمال}
$\frac12$ لكل منهما، على نحو مستقل. ويشتري جامع $5$ علب.
احسب \omterm{def:g10:proba:distribution}{احتمال} أن ينال الجامع تمثالًا واحدًا على الأقل من
كل نوع. (\omterm{def:g10:proba:operations}{المتممة}: كلها A أو كلها B.)
\end{exercise}

\begin{exercise}[$\star\star$]\label{exo:g11:binom:8}
لاعب كرة سلة يسجّل الرميات الحرة \omterm{def:g10:proba:distribution}{باحتمال} $p$،
على نحو مستقل. وليكن $X \sim \mathcal B(3, p)$ عدد التسجيلات في
ثلاث رميات. عبّر عن $\P(X = 3)$ و $\P(X \geq 1)$ بدلالة $p$،
وأوجد من أجل أي $p$ يساوي \omterm{def:g10:proba:distribution}{احتمال} تسجيل الرميات الثلاث
$\frac{27}{64}$.
\end{exercise}

\begin{exercise}[$\star\star$]\label{exo:g11:binom:9}
كم مرة يجب أن تُرمى قطعة نقد غير مغشوشة حتى يتجاوز \omterm{def:g10:proba:distribution}{احتمال} الحصول
على صورة واحدة على الأقل $0.99$؟ (\omterm{def:g10:proba:operations}{المتممة}، ثم جرّب قيمًا
متتابعة للعدد $n$.)
\end{exercise}

\begin{exercise}[$\star\star$]\label{exo:g11:binom:10}
باستعمال قاعدة باسكال (\cref{prop:g11:binom:pascal}) و
$\binom n0 = \binom nn = 1$، برهن على أن مجموع خانات كل سطر من
\omterm{prop:g11:binom:pascal}{مثلث باسكال} هو $2^n$: وفسّر الطرفين على أنهما يعدّان كل
مسارات \omterm{met:g10:proba:tree}{الشجرة}.
\end{exercise}

\begin{exercise}[$\star\star\star$]\label{exo:g11:binom:11}
يدّعي سياسي أن نسبة تأييده $60\%$. وفي \omterm{def:g10:proba:sample}{عينة} عشوائية من $10$ أشخاص،
يؤيده $3$ فقط.
\begin{enumerate}
  \item في ظل هذا الادعاء، أي \omterm{def:g11:prob:rv}{توزيع} يتبعه عدد $X$
        المؤيدين في \omterm{def:g10:proba:sample}{العينة}؟ احسب $\P(X \leq 3)$.
  \item باستعمال قاعدة القرار في \cref{met:g11:binom:decision} بعتبة
        $5\%$، هل الملاحظة متوافقة مع الادعاء؟
\end{enumerate}
\end{exercise}

\section{مسألة: لوح غالتون}

\begin{problem}[{مسألة نهاية الأسبوع --- كرات وأوتاد ومثلث
باسكال: كيف يولد الشكل الجرسي، ولماذا تحابي سلاسل التصفيات
الفريق الأقوى، ومتى نصرخ بالغش}]
\label{pb:g11:binom:1}
أسقط ألف كرة عبر شبكة من الأوتاد، وكل ارتداد
رمية نقد عادلة يمينًا أو يسارًا، فتمتلئ الخانات تحتها في
جرس أملس متناظر --- في كل مرة. وتُسمّى الآلة
لوح غالتون، ورياضياتها هي بالضبط \omterm{def:g11:prob:rv}{التوزيع} الثنائي
في هذا الفصل (\cref{thm:g11:binom:binomial}). وتبني هذه
المسألة المثلث، وتشغّل اللوح، وتحكّم في
سلسلة من سبع مباريات، وتنتهي حيث يكسب \omterm{def:g11:prob:rv}{التوزيع} الثنائي
أجره: في تقرير متى ينبغي لملاحظة أن تجعلنا نشك في
ادعاء.

\medskip
\noindent\textbf{الجزء الأول --- المثلث.}
\begin{enumerate}
  \item ابنِ \omterm{prop:g11:binom:pascal}{مثلث باسكال} حتى السطر $6$
        (\cref{prop:g11:binom:pascal}). وصُغ وفسّر
        التناظر $\binom nk = \binom{n}{n-k}$ في جملة واحدة
        (اختيار $k$ من الأشياء هو نفسه\,\dots).
  \item تحقق على السطرين $4$ و $5$ من أن مجموع كل سطر هو $2^n$،
        وبرهن على ذلك: فماذا تعدّ المقادير $\binom nk$
        مجتمعة؟
  \item استخرج قاعدة باسكال
        $\binom{n+1}{k} = \binom nk + \binom{n}{k-1}$ من جديد بحجة
        اللجنة: ثبّت شخصًا مميزًا واحدًا ثم
        اقسم اللجان حسب مصير ذلك الشخص.
  \item احسب $\binom73$ مرتين: من المثلث، ومن
        صيغة العوامل.
  \item تحقق من متطابقة الدرج
        $\binom22 + \binom32 + \binom42 + \binom52 =
        \binom63$، وفسّرها بتتالي قاعدة باسكال
        نزولًا من $\binom63$.
\end{enumerate}

\medskip
\noindent\textbf{الجزء الثاني --- اللوح.}
تسقط كرة عبر $n$ سطرًا من الأوتاد؛ وعند كل وتد ترتد
يمينًا أو يسارًا \omterm{def:g10:proba:distribution}{باحتمال} $\frac12$، على نحو مستقل. ورقّم
الخانات من $0$ إلى $n$ بعدد الارتدادات إلى اليمين.
\begin{enumerate}[resume]
  \item فسّر، بقائمة التحقق في
        \cref{met:g11:binom:recognize}، لماذا يتبع رقم الخانة
        \omterm{def:g11:prob:rv}{التوزيع} الثنائي
        $\mathcal B\!\left(n, \frac12\right)$.
  \item من أجل لوح صغير ($n = 4$): أعطِ احتمالات الخانات
        الخمس. فأي خانة أكثر ازدحامًا؟
  \item والآن $n = 10$ و $1\,024$ كرة: ما أعداد الكرات
        \emph{المأمولة} في الخانة الوسطى، وفي الخانة
        $7$، وفي كل خانة طرفية؟ صِف شكل الكومة.
  \item من أجل $X \sim \mathcal B\!\left(10, \frac12\right)$:
        احسب $\E(X)$ و $V(X)$ و $\sigma$
        (\cref{prop:g11:binom:expectation})؛ ثم احسب
        نسبة الكرات المأمولة ضمن $2\sigma$ عن
        المركز (الخانات من $2$ إلى $8$) وقارنها بضمانة
        تشيبيشيف في \cref{pb:g11:stat:1}.
  \item لوح مائل يرتد يمينًا \omterm{def:g10:proba:distribution}{باحتمال} $0.6$:
        أعطِ $\E$ و $V$ و $\sigma$ من أجل $n = 10$، وصِف
        ما يحدث للكومة.
  \item في جملة أو جملتين: ما الذي يصنع الشكل الجرسي
        في تصميم اللوح --- ولماذا تتكوّم مقادير كثيرة جدًا
        في العالم الواقعي (الأطوال، وأخطاء القياس)
        بالطريقة نفسها؟ (والمبرهنة العميقة وراء كليهما هي
        مبرهنة النهاية المركزية، قمة درس الاحتمالات في
        الكتب الجامعية.)
\end{enumerate}

\medskip
\noindent\textbf{الجزء الثالث --- الأفضل من سبع.}
فريقان يلعبان سلسلة: الأول الذي يبلغ $4$ انتصارات ينال اللقب؛
والمباريات مستقلة.
\begin{enumerate}[resume]
  \item فريقان متكافئان ($p = \frac12$): احسب \omterm{def:g10:proba:distribution}{احتمال}
        أن تنتهي السلسلة باكتساح ($4$ مباريات بالضبط).
  \item احسب \omterm{def:g10:proba:distribution}{احتمال} أن تبلغ السلسلة المباريات
        $7$ كاملة (ماذا يجب أن تكون النتيجة بعد $6$؟).
  \item أكمل \omterm{def:g11:prob:rv}{توزيع} طول السلسلة
        ($4$ أو $5$ أو $6$ أو $7$ مباريات) من أجل فريقين متكافئين، ثم
        احسب الطول المأمول. فأي الأطوال أكثر
        \omterm{def:g10:proba:distribution}{احتمالًا}؟
  \item والآن يربح فريق كل مباراة \omterm{def:g10:proba:distribution}{باحتمال} $p = 0.6$. احسب
        \omterm{def:g10:proba:distribution}{احتمال} أن ينال السلسلة (بالفوز في $4$ أو $5$ أو $6$
        أو $7$: وفي كل حالة يربح الفريق المباراة الأخيرة و
        $3$ من السابقات). فماذا فعلت السلسلة
        بأفضلية المباراة الواحدة؟
  \item قارن بمباراة نهائية واحدة ($60\,\%$) وبسلسلة من ثلاث
        (احسبها). وصُغ الأثر العام لطول السلسلة
        على المهارة في مواجهة الحظ --- ولماذا تفضّل الدوريات
        النهائيات الطويلة.
\end{enumerate}

\medskip
\noindent\textbf{الجزء الرابع --- متى نصرخ بالغش.}
\begin{enumerate}[resume]
  \item تُرمى قطعة نقد $100$ مرة فتظهر $62$ صورة. من أجل
        قطعة غير مغشوشة، أعطِ $\E$ و $\sigma$ \omterm{pb:g11:stat:1}{والدرجة المعيارية}
        (\cref{pb:g11:stat:1}) للملاحظة. فما الحكم
        باصطلاح $2\sigma$؟
  \item يدّعي مورّد أن نسبة القطع المعيبة $2\,\%$. وفي دفعة
        من $50$ تجد $3$ معيبة. احسب
        $\P(X \geq 3)$ في ظل الادعاء
        ($X \sim \mathcal B(50, 0.02)$؛ ومرّ عبر
        $\P(X = 0), \P(X = 1), \P(X = 2)$). فهل هذا مقلق عند
        عتبة $5\,\%$ (\cref{met:g11:binom:decision} و
        \cref{exo:g11:binom:11})؟
  \item مثابرة اليانصيب: كل تذكرة تربح (شيئًا ما) \omterm{def:g10:proba:distribution}{باحتمال}
        $\frac{1}{1000}$. احسب \omterm{def:g10:proba:distribution}{احتمال}
        فوز واحد على الأقل عند شراء $1\,000$ تذكرة. والجواب
        ($\approx 63\,\%$، لا $100\,\%$!) يخفي ثابتًا شهيرًا:
        احسب $0.999^{1000}$ واحتفظ بالعدد
        $0.368$ في ذهنك للسنة المقبلة.
  \item الخاتمة --- صورة \omterm{def:g11:prob:rv}{التوزيع} الثنائي: قائمة التحقق
        للتعرف عليه (ثبات $n$، والاستقلال، وثبات $p$)؛
        \omterm{prop:g11:binom:pascal}{ومثلث باسكال} جدولًا له؛ والجرس شكلًا له؛
        و $np$ و $np(1-p)$ بوصلة له؛ ووريثاه المنتظران
        في السنة المقبلة --- \omterm{def:g10:functions:graph}{منحنى} الجرس الأملس و
        قانون الأعداد الكبيرة. جملة واحدة لكل بند.
\end{enumerate}
\end{problem}
